\chapter{绘图（Draft）：2D 线条与工作平面}

Draft 工作台适合画二维图：平面上的线、圆、多边形、标注，以及路径、阵列这类“图形式”的操作。
它的坐标系是\textbf{工作平面}，不是固定的 XY 平面——这是用好 Draft 的关键。

\section{入口与工具条}

从 \menupath{更多 ▾ → Draft} 进入，或使用 \ui{零件} 页里的 \ui{Draft} 工具组。工具按 FreeCAD 的
五条工具条分组：\ui{创建}、\ui{标注}、\ui{修改}、\ui{实用工具}、\ui{捕捉}。

\section{工作平面}

在 \ui{更多工具} 里选择 \ui{工作平面}，有五种：

\begin{center}
\small
\begin{tabularx}{\linewidth}{@{}lX@{}}
\toprule
\textbf{选项} & \textbf{含义} \\
\midrule
上 (XY) / 前 (XZ) / 侧 (YZ) & 三个坐标平面 \\
所选面 & 用当前选中的平面作为工作平面（先选中，否则提示“请先选择平面”；非平面会报错）\\
当前视图 & 平面过相机目标点，屏幕右/上为 u/v，视线为法向 \\
\bottomrule
\end{tabularx}
\end{center}

工作平面影响落点投影、矩形/圆/多边形的朝向、镜像与旋转轴、键入坐标的解释，以及栅格捕捉的取整方向。
它\textbf{不影响} 3D 网格的显示（网格始终在 XY 地面）。

\section{捕捉}

Draft 的捕捉条有 15 个开关，\ui{Lock} 是总开关：

\begin{itemize}
  \item \textbf{默认开启}：Lock、端点、中点、圆心、近点、交点、垂足、栅格、工作平面；
  \item \textbf{默认关闭}：角度、正交、平行、延伸、特殊、尺寸；
  \item 端点/中点/圆心取自可见形状：顶点、每条边的中点、圆/弧的中心，取 18 dp 内最近者；
  \item 没有对象可捕捉时，落点落在工作平面上：\ui{正交}锁定水平/竖直，否则 \ui{角度} 每 15° 吸附，
        最后 \ui{栅格} 按 1 mm 取整；
  \item 捕捉名称会显示在光标旁（例如“端点”“交点”）。
\end{itemize}

\section{点输入}

轻触落点；顶部横幅显示工具名、提示，以及 \ui{Typing}/\ui{撤销}/\ui{关闭}/\ui{完成}/\ui{复制}/✕。

\begin{itemize}
  \item \ui{Type}：输入 \file{x, y, z} 或 \file{@dx, dy, dz}（2 或 3 个数，逗号/分号/空格分隔）。
        坐标按\textbf{工作平面的局部坐标}解释，\file{@} 表示相对上一点；
  \item 长按进入精确模式（准星在手指上方 60 dp），抬起即落点；
  \item 多段工具双击结束（第二次轻触会撤掉刚加的点），定长工具画够就自动结束；
  \item \ui{完成} 需要达到最少点数；\ui{关闭} 只对折线且至少 2 点；\ui{撤销} 删掉最后一点。
\end{itemize}

\section{绘制工具}

\begin{longtable}{@{}p{0.18\linewidth}p{0.10\linewidth}p{0.62\linewidth}@{}}
\toprule \textbf{工具} & \textbf{点数} & \textbf{说明} \\ \midrule \endhead
线 & 2 & 起点 → 终点 \\
折线 & 多点 & 连点；双击结束 \\
矩形 & 2 & 对角两点 \\
圆 & 2 & 圆心 → 半径点 \\
圆弧 & 3 & 圆心 → 起点 → 终点 \\
三点圆弧 & 3 & 弧上三点 \\
椭圆 & 2 & 对角两点（长短轴自动交换）\\
多边形 & 2 & 先问边数（默认 6），再圆心 + 顶点（内接）\\
B 样条 & 多点 & 连点 \\
贝塞尔 & 多点 & 三次曲线，每段 4 点后继续 \\
点 & 1 & 单点 \\
文字 & 2 & 先输入文字（\file{\textbackslash n} 换行），再点位置 \\
形状字符串 & 2 & 再问高度（默认 10）；字体固定为 Roboto \\
标注 & 3 & 三点线性标注 \\
标签 & 多点 & 先输文字；目标取当前选择；再点目标点与标签点 \\
面 & — & 先选面再执行（面绑定）\\
填充 & — & 先选对象；选 8 种图案之一与比例（默认 10）\\
\bottomrule
\end{longtable}

\section{修改与阵列}

使用前\textbf{先选对象}，否则提示“请先选择对象”。

\begin{center}
\small
\begin{tabularx}{\linewidth}{@{}lX@{}}
\toprule
\textbf{工具} & \textbf{操作} \\
\midrule
移动 & 基点 → 目标点（横幅的“复制”开启则复制）\\
旋转 & 圆心 → 基准角点 → 旋转角点（轴为平面法向）\\
缩放 & 基点 → 输入因子（默认 2）\\
镜像 & 两点定镜像线 \\
偏移 & 点一下定侧与距离；圆/弧改半径，线按垂距；只处理第一个对象 \\
拉伸 & 框两点 → 基点 → 目标点 \\
修剪 / 延伸 & 点目标（面则沿法向拉伸，线则改最近端）\\
分割 & 先选线，再点分割位置 \\
编辑 & 点节点，再点新位置（仅对带 Points 的对象）\\
\bottomrule
\end{tabularx}
\end{center}

其它：\ui{升级}/\ui{降级}（会删原对象）、\ui{连接}、\ui{Draft↔Sketch 转换}、\ui{克隆}、
\ui{2D 形状视图}（固定 +Z）、\ui{线转 B 样条}、\ui{修复}、\ui{加入施工组}。

\subsection{阵列}

\begin{itemize}
  \item \ui{矩形阵列}：问 X/Y 份数与间隔（默认 2/2/100；Z 方向 1 且间隔与 X 相同）；
  \item \ui{极坐标阵列}：问元素数（5）与扫掠角（360°），\textbf{中心固定在原点}，不能点选；
  \item \ui{环形阵列}：问径向 100、切向 50、环数 3，中心同样在原点；
  \item \ui{路径阵列 / 扭曲路径阵列 / 点阵列}：按选择顺序（基体在前，路径或点对象在后），
        数量固定（4 或 15），没有参数对话框；
  \item Link 版阵列为链接阵列（引用原对象）。
\end{itemize}

\section{样式与图层}

\begin{itemize}
  \item \ui{设置样式}/\ui{应用当前样式}/\ui{标注样式编辑器} 都转到标准的“外观”对话框：
        \textbf{只有 12 色 + Default}，用于改对象颜色；移动端没有线宽、线型、图层样式与注释样式编辑器；
  \item 图层：\ui{新建图层}（有选择则把选择放入）、\ui{添加到图层}、\ui{图层管理器}（点图层切换显示）、
        \ui{自动分组}（设置新对象的默认图层）、\ui{添加到组}/\ui{选择组}；
  \item 施工模式：把新对象加入 Construction 组。
\end{itemize}

\begin{practice}[练习：带标注的平面布置图]
\begin{enumerate}
  \item 工作平面选 \ui{上 (XY)}，打开栅格捕捉；
  \item 用 \ui{矩形} 画 4000 × 3000 的房间（Type 输入对角坐标），再用 \ui{折线} 画一条隔墙；
  \item 用 \ui{标注} 标出两道主要尺寸；
  \item 用 \ui{矩形阵列} 复制一列柱网（X 间隔 1000、2 份）；
  \item 保存，转到 \ui{更多 → 零件} 用 \ui{拉伸} 把墙线变成 3D 墙体。
\end{enumerate}
\end{practice}

\section{源码可见的限制}

\begin{itemize}
  \item 样式只有颜色；极坐标/环形阵列的中心固定在原点；矩形阵列的 XYZ 共用一个间隔；
  \item 路径/点阵列没有参数对话框，且依赖选择顺序；
  \item 平行、延伸、特殊、尺寸这四个捕捉开关在源码中没有对应逻辑；
  \item 栅格间隔固定 1 mm，且网格画在 XY 地面而不是工作平面；
  \item 偏移只处理第一个对象；修剪/延伸只支持线/面；编辑只对有 Points 的对象；
  \item 形状字符串字体固定 Roboto；2D 形状视图固定 +Z；填充作用于整个对象。
\end{itemize}
